\documentclass{article}

\usepackage{makecell}

\title{Ch. 22 CW \#1}
\author{Principles of Microeconomics}
\date{\today}

\begin{document}

\maketitle

\subsection*{Instructions}

Please justify all of your answers and submit your solutions on Gradescope by the beginning of next class.

\subsection*{Problems}

\begin{enumerate}

\item Raj consumes only cheese and crackers.
	
	\begin{enumerate}
	
	\item Could cheese and crackers both be inferior goods for Raj? Explain
	
	\item Suppose that cheese is a normal good for Raj and crackers are an inferior good. If the price of cheese falls, what happens to Raj's consumption of crackers? What happens to his consumption of cheese? Explain.
	
	\end{enumerate}

\item Darius buys only milk and cookies.

	\begin{enumerate}
	
	\item In year 1, Darius earns \$100, milk costs \$2 per quart, and cookies cost \$4 per dozen. Draw Darius's budget constraint.
	
	\item Now suppose that all prices increase by 10\% in year 2 and that Darius's salary increases by 10 percernt as well. Draw Darius's new budget constraint. How would Darius's optimal bundle of milk and cookies in year 2 compare with his optimal bundle in year 1?
	
	\end{enumerate}
	
\item The economist George Stigler once wrote that, according to consumer theory,``if consumers do not buy less of a commodity when their incomes rise, they will surely buy less when the price of the commodity rises." Explain this statement using the concepts of income and substitution effects.

\newpage

\item Five consumers have the following marginal utility of apples and pears:

	\begin{center}
	\begin{tabular}{| c | c | c |}
	\hline
	& \makecell{Marginal Utility of\\Apples} & \makecell{Marginal Utility of\\Pears} \\
	\hline
	Claire & 6 & 12 \\
	\hline
	Phil & 6 & 6 \\
	\hline
	Haley & 6 & 3 \\
	\hline
	Alex & 3 & 6 \\
	\hline
	Luke & 3 & 12 \\
	\hline
	\end{tabular}
	\end{center}
	
	The price of an apple is \$1, and the price of a pear is \$2. Which, if any, of these consumers are optimizing their choices of fruit? For those who are not, how should they change their spending?
	
\item Draw the indifference curve for someone deciding how to allocate time between work and leisure. Suppose the wage increases. Is it possible that the person's consumption would fall? Is this plausible? Discuss. (Hint: Think about income and substitution effects.)
		
\end{enumerate}

\end{document}